\myChapter{Revue de la littérature}{Revue de la littérature}
\label{chap:revue-litterature}
\myMiniToc{mysection}{Sommaire}
\nocite{prisma2020,schaffer2021,sajji2022,sajji2023,pietruck2020,greifenberg2015,rahad2021,dragas2024,delgado2022,avila2021,lopes2023,bucaioni2022,diruscio2022,rokis2023,ajimati2025,castillo2024,sajji2025,krouwel2024,martinez2026,drave2022,fedeli2024,soliman2025,horner2026,kourani2025,kourani2024,licardo2026,neto2022,uyanik2023,hoyos2017,hemel2008,syriani2018,siahaan2024,laaz2020,salado2024,campagna2025,pankowska2024,diaz2018}


\mySection{Introduction}{}

La transformation de modèles de processus métier en applications web constitue un défi scientifique et technique qui mobilise plusieurs champs disciplinaires : la modélisation par BPMN, les systèmes de gestion de processus, l'ingénierie dirigée par les modèles et les plateformes low-code. Si chacun de ces domaines a fait l'objet de contributions significatives, les travaux qui les articulent de manière cohérente en vue de produire des applications web complètes, maintenables et dont le code source reste maîtrisable, demeurent rares et fragmentés.

Ce chapitre propose une revue systématique de la littérature visant à analyser les travaux existants dans ce paysage scientifique. La démarche, conduite sur un corpus de travaux publiés entre 2015 et 2026, s'inspire des principes de rigueur et de transparence recommandés par PRISMA 2020, adaptés au contexte de l'informatique et du génie logiciel. Elle est guidée par une question centrale : quelles approches permettent de passer de modèles BPMN vers des artefacts logiciels web, et quelles limites restent ouvertes pour une recherche portant sur la génération d'applications web à partir de processus métier ?

L'analyse est conduite selon quatre axes. Le premier porte sur BPMN en tant que langage de modélisation et support à l'automatisation des systèmes orientés processus. Le deuxième examine les BPMS classiques, leurs capacités de déploiement applicatif et leurs limites structurelles face aux exigences d'une application web complète. Le troisième traite de l'ingénierie dirigée par les modèles et des mécanismes de génération de code, en interrogeant leur aptitude à produire des artefacts cohérents, évolutifs et intégrables. Le quatrième confronte les approches model-driven aux plateformes low-code et no-code, dont l'essor récent soulève des questions de gouvernance, de maintenabilité et de dépendance technologique.

\mySection{Question de recherche et grille d'analyse}{}

\subsection{Questions de recherche et motivations}

La question principale retenue est la suivante :

\begin{quote}
\textbf{RQ0 : Quelles approches permettent de passer de modèles BPMN ou de modèles de processus métier vers des applications web ou des artefacts logiciels, et quelles limites restent ouvertes dans les travaux existants ?}
\end{quote}

Cette question permet d'orienter la revue vers les travaux qui ne se limitent pas à décrire BPMN comme notation, mais qui cherchent à l'utiliser comme support de conception, d'automatisation, de génération ou d'intégration logicielle. L'objectif est d'identifier ce que les approches existantes produisent réellement, ce qu'elles automatisent, ce qu'elles laissent à la charge du développeur et les difficultés qu'elles ne résolvent pas encore.

La question principale est déclinée en six questions secondaires. Pour chacune d'elles, une motivation précise est formulée afin de clarifier son rôle dans l'analyse.

\begin{table}[H]
\caption{Questions secondaires et motivations de la revue de littérature}
\label{tab:questions-recherche-motivations}
\centering
\footnotesize
\renewcommand{\arraystretch}{1.2}
\begin{tabular}{|p{1.3cm}|p{6.0cm}|p{6.0cm}|}
\hline
\textbf{Code} & \textbf{Question secondaire} & \textbf{Motivation} \\
\hline
QR1 & Quel rôle BPMN joue-t-il dans la conception et l'automatisation des systèmes d'information orientés processus ? & BPMN est souvent présenté comme un langage de communication entre acteurs métier et acteurs techniques. Il est donc nécessaire de déterminer s'il reste un simple support documentaire ou s'il peut être exploité dans une chaîne de transformation logicielle. \\
\hline
QR2 & Quelles sont les capacités et les limites des BPMS classiques pour construire des applications web métier complètes et maîtrisables ? & Les BPMS exécutent et supervisent des processus, mais une application web complète implique aussi des interfaces, des données, des règles de sécurité, des API, des tests et du code maintenable. Cette question permet de distinguer l'exécution de workflow de la construction d'une application logicielle complète. \\
\hline
QR3 & Comment les travaux existants transforment-ils BPMN ou des modèles de processus en interfaces, API, services ou code applicatif ? & La transformation de modèles vers des artefacts web constitue le coeur du problème étudié. Il faut donc analyser les types d'artefacts produits, le niveau d'automatisation atteint et les informations complémentaires nécessaires à la génération. \\
\hline
QR4 & Comment l'ingénierie dirigée par les modèles et la génération de code traitent-elles la qualité, la traçabilité et l'évolution des artefacts produits ? & Générer du code ne suffit pas ; le code produit doit rester compréhensible, modifiable, versionnable et compatible avec l'évolution du modèle. Cette question permet d'étudier les mécanismes d'intégration entre modèles, code généré et code manuel. \\
\hline
QR5 & Comment les approches low-code, no-code et model-driven se rapprochent-elles ou se différencient-elles ? & Les plateformes low-code promettent une accélération du développement à partir de modèles visuels. Il est donc utile de comparer leur logique à celle de l'ingénierie dirigée par les modèles, notamment sur la maîtrise du code, la dépendance aux plateformes et la gouvernance technique. \\
\hline
QR6 & Quelles lacunes scientifiques et techniques persistent dans les travaux récents ? & Cette question permet de synthétiser les limites observées : couverture partielle de la génération, insuffisance de BPMN seul, faible prise en compte de la régénération, difficulté à préserver les personnalisations et manque de modularité. \\
\hline
\end{tabular}
\end{table}

\subsection{Grille d'analyse}

Les travaux retenus sont analysés à partir de cinq critères principaux :

\begin{itemize}
    \item \textbf{le rôle du modèle d'entrée :} le modèle est-il uniquement documentaire, exécutable ou transformable ?
    \item \textbf{les artefacts produits :} l'approche génère-t-elle des interfaces, des API, des services, des tests, une base de données ou uniquement des configurations ?
    \item \textbf{le niveau d'automatisation :} la transformation est-elle manuelle, semi-automatique ou automatique ?
    \item \textbf{la maîtrise du code :} le code produit est-il lisible, versionnable, modifiable et maintenable ?
    \item \textbf{l'évolution :} l'approche prévoit-elle la régénération, la préservation des personnalisations et la gestion des conflits ?
\end{itemize}

Cette grille permet d'éviter une simple juxtaposition de résumés d'articles. Elle favorise une lecture critique des contributions et met progressivement en évidence les lacunes utiles au positionnement de notre recherche.

\mySection{BPMN dans la conception des systèmes d'information orientés processus}{}

\subsection{BPMN comme langage de médiation entre métier et technique}

BPMN occupe une place importante dans la conception des systèmes d'information orientés processus, car il permet de représenter graphiquement l'enchaînement des activités, des décisions, des événements et des responsabilités. Son intérêt ne réside pas uniquement dans sa capacité descriptive. Il tient également à son rôle de médiation entre les acteurs métier, qui expriment les besoins organisationnels, et les acteurs techniques, qui doivent traduire ces besoins en structures logicielles.

Schäffer et al. (2021) montrent que BPMN peut servir de base à une plateforme web orientée processus lorsqu'il est intégré à une architecture logicielle appropriée et à des moteurs de processus [2]. Leur contribution illustre la possibilité d'utiliser BPMN comme support de structuration du comportement métier dans une application web. L'intérêt d'une telle approche est de maintenir un lien fort entre la représentation du processus et son exécution informatique. Toutefois, cette perspective reste principalement centrée sur l'orchestration et la supervision des processus, plutôt que sur la génération complète d'un projet web maîtrisable dans son code source.

L'utilisation de BPMN comme point d'entrée présente donc un avantage méthodologique important : elle permet de partir d'un modèle compréhensible par les acteurs métier. Dans un contexte de développement d'applications web, cette caractéristique est essentielle, car elle réduit l'écart entre la compréhension organisationnelle du problème et sa traduction technique. Cependant, ce rapprochement ne suffit pas à lui seul. Le passage du modèle au logiciel exige des transformations, des enrichissements et des décisions d'architecture.

\subsection{Qualité et lisibilité des modèles de processus}

Un modèle BPMN ne peut servir de base à une génération fiable que s'il est suffisamment clair, cohérent et exploitable. La qualité du modèle devient alors un enjeu central. Avila et al. (2021) réalisent une revue systématique consacrée aux lignes directrices de modélisation des processus et à leur support empirique [10]. Leur travail montre que la lisibilité, la cohérence, la taille du modèle, la structuration des activités et la réduction de la complexité influencent la compréhension des modèles de processus.

Dans une perspective de génération d'applications, ces observations sont particulièrement importantes. Un modèle BPMN ambigu ou mal structuré peut conduire à des artefacts logiciels incohérents. Par exemple, une activité mal nommée peut entraîner la génération de composants peu compréhensibles ; une passerelle mal définie peut produire une logique conditionnelle incorrecte ; une absence de distinction entre tâches humaines et tâches automatisées peut rendre difficile la génération des formulaires, services ou contrôleurs appropriés.

Lopes et Guerreiro (2023) insistent également sur l'importance du test et de la vérification des modèles BPMN [11]. Leur revue montre que les techniques de vérification formelle et de test peuvent contribuer à détecter des erreurs dans les modèles de processus avant leur exécution ou leur transformation. Dans le cadre de notre recherche, cette idée est déterminante : la génération d'une application web ne devrait pas commencer uniquement par la lecture syntaxique d'un fichier BPMN, mais par une étape de validation du modèle et de contrôle de sa cohérence.

Les travaux récents sur l'assistance à la modélisation par grands modèles de langage renforcent cette attention portée à la qualité du modèle d'entrée. Kourani et al. proposent d'utiliser les LLM pour générer et raffiner des modèles de processus à partir de descriptions textuelles, avec des mécanismes de contrôle et de correction des erreurs \cite{kourani2024}. Dans une perspective proche, Hörner et al. étudient la génération automatique de modèles BPMN 2.0 depuis des descriptions en langage naturel, en insistant sur les difficultés de qualité et d'évaluation \cite{horner2026}. Ces travaux ne remplacent pas les méthodes classiques de modélisation, mais ils montrent que la production du modèle BPMN elle-même devient un objet de recherche. Pour notre recherche, cela signifie que la chaîne de génération doit rester indépendante de la manière dont le BPMN est produit : que le modèle soit dessiné manuellement ou assisté par l'IA, il doit être vérifié avant d'être transformé.

\subsection{Limites de BPMN pour la génération d'applications web}

Malgré son intérêt, BPMN reste insuffisant pour décrire l'ensemble des informations nécessaires à la construction d'une application web complète. Il permet de représenter les activités et les transitions, mais il ne spécifie pas de manière détaillée les champs des formulaires, les règles de validation, les composants d'interface, les entités persistantes, les politiques d'autorisation ou les scripts de déploiement.

Cette limite explique pourquoi plusieurs travaux associent BPMN à d'autres modèles. Sajji et al. (2022) proposent une transformation semi-automatique de BPMN vers IFML dans une logique MDA [3]. L'approche consiste à partir d'un modèle de processus au niveau CIM, puis à produire un modèle d'interaction au niveau PIM. Cette contribution montre que BPMN peut être utilisé comme point de départ pour dériver des interfaces web, mais elle confirme également qu'un modèle complémentaire est nécessaire pour représenter les aspects interactionnels de l'application.

Sajji et al. (2023) prolongent cette perspective en étudiant la génération d'interfaces graphiques à partir de BPMN via IFML jusqu'au niveau PSM [4]. Leur travail met en évidence la possibilité de rapprocher davantage le processus métier de l'interface utilisateur. Toutefois, l'approche reste principalement focalisée sur la partie interactionnelle. Elle ne couvre pas à elle seule l'ensemble d'une architecture applicative comprenant les services backend, les règles de sécurité, les tests, la gestion des rôles et la régénération contrôlée du code.

Ainsi, BPMN apparaît comme un excellent point d'entrée métier, mais non comme un modèle complet d'application. Pour une recherche portant sur la génération d'applications web à partir de BPMN, cette conclusion conduit à deux exigences. Premièrement, BPMN doit être enrichi ou interprété par des contrats intermédiaires. Deuxièmement, la chaîne de génération doit distinguer clairement ce qui relève du processus, de l'interface, de la donnée, du service, de la sécurité et de la configuration technique.

\mySection{BPMS classiques et applications web métier}{}

\subsection{Apports des BPMS}

Les BPMS, ou \textit{Business Process Management Systems}, constituent une famille d'outils permettant de modéliser, exécuter, suivre et parfois optimiser des processus métier. Leur contribution principale réside dans leur capacité à gérer des instances de processus, à affecter des tâches, à contrôler des transitions, à intégrer des événements et à fournir des mécanismes de supervision. Ils apportent donc une réponse concrète au besoin d'automatisation des workflows.

Dans les architectures orientées processus, les BPMS peuvent jouer le rôle de moteur d'exécution. Ils prennent en charge la logique de contrôle du processus et permettent aux organisations de suivre l'état des dossiers ou des demandes. Schäffer et al. (2021) montrent que cette logique peut être exploitée dans une plateforme web orientée processus, à condition d'être intégrée à une architecture de référence claire [2]. Cette approche confirme l'intérêt des moteurs de processus pour structurer l'exécution des activités métier.

Cependant, la finalité d'un BPMS n'est pas nécessairement de produire une application web complète sous forme d'un projet logiciel autonome. Un BPMS peut orchestrer un processus, mais l'application web métier exige généralement d'autres éléments : modèles de données, interfaces spécialisées, logique métier propre au domaine, API, mécanismes de sécurité, tests, déploiement et maintenance du code. Cette distinction est essentielle pour comprendre les limites des BPMS classiques dans le contexte de notre recherche.

\subsection{Écart entre moteur de workflow et application web complète}

Le passage d'un processus exécutable à une application web complète reste un problème complexe. Une application métier ne se limite pas à l'exécution d'une suite d'activités. Elle doit proposer des écrans adaptés aux rôles des utilisateurs, stocker des données persistantes, appliquer des règles de validation, exposer des services, contrôler les droits d'accès et permettre l'évolution fonctionnelle du système.

Pietruck et Steffens (2020) abordent cette difficulté en proposant la génération d'API REST pour des moteurs de workflow à partir de modèles de processus et de données [5]. Leur contribution est intéressante, car elle montre que les API fournies par les moteurs de processus ne correspondent pas toujours directement aux concepts métier manipulés par les applications clientes. Les auteurs cherchent donc à produire une interface plus orientée domaine, capable de masquer une partie de la complexité du moteur de workflow. Cette idée rejoint fortement les préoccupations de notre recherche, notamment la nécessité de produire des contrats et services compréhensibles par les développeurs d'applications web.

Toutefois, la génération d'API autour d'un moteur de workflow ne suffit pas à résoudre tout le problème de la génération d'application. Elle améliore l'intégration entre processus et composants logiciels, mais ne produit pas nécessairement un frontend complet, une architecture backend modulaire, des tests applicatifs, des politiques de sécurité ou un mécanisme de régénération contrôlée. Cette limite montre que l'approche par moteur de workflow doit être complétée par une approche de génération logicielle plus large.

\subsection{Limites des BPMS pour une approche source-owned}

Les BPMS classiques répondent principalement à une logique d'exécution, de supervision et d'orchestration. Ils sont utiles lorsque l'objectif est de piloter des processus dans un environnement centralisé. En revanche, ils peuvent être moins adaptés lorsqu'une équipe souhaite disposer d'un projet web complet, versionnable, personnalisable et maintenable comme n'importe quelle application développée manuellement.

La notion de \textit{source-owned} désigne ici une orientation selon laquelle le code généré reste sous la maîtrise de l'équipe de développement. Cela signifie que l'application produite doit pouvoir être consultée, modifiée, testée, versionnée et déployée sans dépendre entièrement d'une plateforme fermée. Cette exigence ne revient pas à rejeter les moteurs de processus, mais à refuser que l'architecture de l'application soit entièrement prisonnière d'un moteur ou d'un outil externe.

Dans cette perspective, les travaux orientés génération se distinguent d'un BPMS classique. Ils ne visent pas uniquement à exécuter des workflows, mais à produire les éléments constitutifs d'une application web métier à partir de modèles BPMN enrichis. Cette orientation suppose de transformer le modèle de processus en artefacts exploitables : routes, contrôleurs, services, vues, formulaires, contrats, règles d'accès, tests et configurations. Le moteur de processus peut éventuellement rester une option, mais il ne doit pas nécessairement être le centre unique de l'application générée.

\subsection{Applications web orientées processus}

Les systèmes d'information orientés processus, ou \textit{process-aware information systems}, constituent une famille de travaux particulièrement proche de notre problématique. Drave et al. (2022) montrent que l'ingénierie dirigée par les modèles peut être utilisée pour concevoir de tels systèmes en reliant modèles de processus, artefacts logiciels et mécanismes d'exécution \cite{drave2022}. Cette approche confirme que le processus métier peut être plus qu'une documentation : il peut devenir une source structurante pour l'application.

Neto et al. (2022) proposent avec AKIP une plateforme de développement d'applications web orientées processus \cite{neto2022}. Leur travail est important parce qu'il rapproche explicitement la logique BPMN de la production d'applications web. Il montre que les organisations ont besoin d'outils capables de produire rapidement des applications centrées sur les tâches, les rôles et le suivi des instances de processus. Toutefois, l'approche reste fortement liée à une plateforme d'automatisation. La question de la possession complète du code source, de sa modification hors plateforme et de sa régénération contrôlée reste donc ouverte.

Uyanik et Sayar (2023) abordent également le développement web orienté processus avec génération automatique de code \cite{uyanik2023}. Leur contribution va dans le sens d'une automatisation plus poussée, mais elle illustre aussi une difficulté récurrente : générer une application opérationnelle exige de coordonner plusieurs niveaux, notamment le processus, les formulaires, les données, les services et l'exécution. Ces travaux confortent l'idée qu'une application web métier ne peut pas être déduite d'un diagramme BPMN par simple correspondance élément par élément.

Enfin, Fedeli et al. (2024), avec FloBP, montrent que BPMN peut être étendu dans une approche model-driven pour développer et exécuter des processus métier enrichis par des objets IoT \cite{fedeli2024}. Même si leur domaine d'application est différent, leur contribution rappelle qu'un processus métier devient souvent insuffisant lorsqu'il doit dialoguer avec un environnement technique spécialisé. Cette observation rejoint directement le besoin d'artefacts intermédiaires et de modèles complémentaires.

\mySection{Transformation de modèles BPMN vers des artefacts web}{}

\subsection{Transformation BPMN vers modèles d'interaction}

Les travaux de transformation BPMN vers des modèles d'interaction montrent que le processus métier peut servir de base à la conception d'interfaces. Sajji et al. (2022) proposent une méthodologie de transformation BPMN vers IFML dans une approche MDA [3]. Leur démarche part d'un modèle de niveau CIM, représentant les processus métier, pour produire un modèle de niveau PIM centré sur les interactions utilisateur. Ce type de transformation permet de rapprocher les activités humaines du processus des écrans ou vues qui doivent être proposés aux utilisateurs.

Cette contribution est importante pour notre recherche, car elle confirme que BPMN peut être interprété comme une source d'information pour la génération d'interfaces. Une tâche utilisateur dans BPMN peut, par exemple, correspondre à un formulaire de saisie, une page de validation ou une action métier. Une passerelle peut entraîner des choix conditionnels dans l'interface. Un événement peut être associé à une notification ou à une transition automatique. Toutefois, cette correspondance n'est pas toujours directe. Elle exige des règles de transformation précises et parfois des informations complémentaires non présentes dans le modèle BPMN initial.

Sajji et al. (2023) vont plus loin en étudiant la génération d'interfaces graphiques à partir de BPMN via IFML jusqu'au niveau PSM [4]. Cette extension permet de rapprocher la modélisation abstraite des artefacts techniques propres à une plateforme. Elle constitue une avancée importante, car elle montre que la chaîne de transformation peut être prolongée vers des éléments plus concrets. Néanmoins, l'approche reste centrée sur l'interface et ne traite pas complètement les aspects backend, sécurité, tests, intégration continue ou gestion de l'évolution du code.

Laaz et al. (2020) enrichissent cette famille de travaux en combinant des ontologies de domaine et des modèles BPMN au niveau CIM afin de générer des modèles IFML \cite{laaz2020}. L'intérêt d'une telle approche est de reconnaître que BPMN ne porte pas toujours seul la connaissance nécessaire à la génération d'interfaces. Les concepts du domaine, les relations entre entités et le vocabulaire métier peuvent aider à produire des modèles d'interaction plus cohérents. Cette idée est importante pour notre recherche, car elle confirme que le modèle BPMN doit parfois être complété par une couche sémantique ou contractuelle.

Díaz et al. (2018) proposent, de leur côté, une méthode de génération de prototypes d'interfaces à partir de BPMN et de diagrammes de classes \cite{diaz2018}. Leur travail identifie des correspondances entre certains patrons BPMN et des composants d'interface, comme les formulaires, les assistants par étapes, les groupes de champs ou les choix conditionnels. Cette contribution a une valeur pratique, car elle montre comment des tâches utilisateur, des passerelles exclusives ou des synchronisations peuvent inspirer des écrans concrets. Elle révèle cependant une limite importante : lorsque la génération dépend fortement de stéréotypes ou de diagrammes complémentaires, elle devient moins directement applicable à des BPMN ordinaires produits par des acteurs métier.

Les travaux plus récents de Sajji et al. (2025), qui étudient le passage d'IFML vers Symfony, prolongent cette logique de transformation vers un framework web concret \cite{sajji2025}. Ils montrent que la génération devient plus réaliste lorsqu'elle s'appuie sur une chaîne de modèles successifs plutôt que sur une traduction immédiate du processus vers le code. Cette observation renforce l'intérêt d'une architecture de génération progressive, dans laquelle BPMN alimente des représentations intermédiaires avant la production des artefacts techniques.

\subsection{Transformation vers API et services backend}

Une application web moderne repose généralement sur une séparation entre frontend, backend, données et services. Dans ce contexte, la transformation d'un modèle de processus ne doit pas seulement produire des écrans. Elle doit également produire des API et des services capables de porter la logique métier. Pietruck et Steffens (2020) proposent précisément une approche de génération d'API REST pour des moteurs de workflow à partir de modèles de processus et de données [5].

Leur travail met en évidence un problème fréquent : les moteurs de workflow exposent souvent des API techniques centrées sur les concepts du moteur, comme les instances, les tâches, les variables ou les transitions. Or, les développeurs d'applications web manipulent plutôt des concepts métier : demande, dossier, validation, acteur, décision, document, notification. La génération d'une API orientée domaine permet donc de réduire l'écart entre le moteur de processus et l'application cliente.

Pour la problématique étudiée, cette idée est fondamentale. Une approche de génération ne doit pas se limiter à traduire les éléments BPMN en fichiers techniques. Elle doit produire des contrats intermédiaires qui expriment les concepts applicatifs de manière stable. Ces contrats peuvent ensuite servir à générer des routes backend, des services, des contrôleurs, des validations et des interfaces. La génération devient ainsi plus maîtrisable, car elle ne dépend pas directement du modèle BPMN brut, mais d'une représentation intermédiaire plus proche de l'architecture logicielle cible.

Dans la même direction, Hoyos et Restrepo-Calle (2017) étudient la génération automatique de code source pour des systèmes d'information web orientés processus \cite{hoyos2017}. Leur contribution confirme que la génération d'applications web ne se limite pas à produire une couche de présentation : elle doit aussi prendre en charge la logique de navigation, les opérations métier et la persistance. Elle met également en évidence une tension classique des approches génératives : plus le générateur produit d'artefacts, plus il doit expliciter ses hypothèses d'architecture afin que le code obtenu reste compréhensible.

\subsection{Nécessité d'artefacts intermédiaires}

La transformation directe d'un modèle BPMN vers du code applicatif présente plusieurs risques. Elle peut créer un couplage trop fort entre la notation graphique et l'architecture cible. Elle peut aussi rendre les règles de génération difficiles à maintenir, car chaque changement dans le modèle ou dans la plateforme cible risque d'affecter directement les fichiers produits. Pour limiter ces difficultés, une approche plus robuste consiste à introduire des artefacts intermédiaires.

Ces artefacts peuvent prendre la forme de contrats décrivant les entités, les transitions, les rôles, les écrans, les actions, les validations et les dépendances fonctionnelles. Ils jouent un rôle de médiation entre le modèle BPMN et les générateurs de code. Cette idée rejoint les principes de l'ingénierie dirigée par les modèles, dans laquelle les transformations successives permettent de passer d'un niveau d'abstraction à un autre.

Dans cette perspective, les contrats intermédiaires peuvent remplir trois fonctions. Premièrement, ils stabilisent l'information extraite du modèle BPMN. Deuxièmement, ils permettent de vérifier la cohérence avant génération. Troisièmement, ils facilitent la séparation entre les générateurs backend, frontend, sécurité, tests et déploiement. Cette séparation est essentielle pour obtenir une architecture modulaire.

Siahaan et al. (2024), dans leur revue sur la génération automatique de transformations entre artefacts de développement logiciel, montrent que le passage d'un artefact à un autre soulève des questions de cohérence, de couverture, de validation et de traçabilité \cite{siahaan2024}. Cette observation dépasse le seul cas de BPMN. Elle rappelle qu'une chaîne de génération robuste doit rendre explicites ses transformations, ses entrées, ses sorties et ses limites. Cette exigence plaide en faveur d'une architecture où chaque étape de transformation peut être inspectée et testée.

\mySection{Ingénierie dirigée par les modèles et génération de code}{}

\subsection{Apport de l'ingénierie dirigée par les modèles}

L'ingénierie dirigée par les modèles, ou MDE pour \textit{Model-Driven Engineering}, repose sur l'idée que les modèles peuvent devenir des artefacts centraux du développement logiciel. Ils ne servent pas seulement à documenter le système, mais peuvent être transformés, validés et utilisés pour générer d'autres artefacts. Dans le contexte de ce mémoire, cette approche fournit un cadre théorique important pour comprendre le passage de BPMN vers une application web.

Di Ruscio et al. (2022) montrent que les approches low-code et l'ingénierie dirigée par les modèles partagent plusieurs objectifs : réduire la complexité du développement, augmenter le niveau d'abstraction et automatiser certaines tâches répétitives [13]. Cette analyse est utile, car elle situe la problématique étudiée dans une continuité scientifique plus large. Il ne s'agit pas seulement d'un problème de génération technique, mais d'une question de transformation de modèles vers des artefacts logiciels.

Castillo et Dávila (2024) proposent une étude de cartographie systématique des outils de développement web dirigé par les modèles [16]. Leur travail confirme que le développement web model-driven couvre plusieurs dimensions : contenu, navigation, présentation, logique métier, données et transformation de modèles. Cette diversité montre que la génération d'applications web n'est pas un problème unique, mais un ensemble de sous-problèmes interdépendants.

\subsection{Génération de code et transformation progressive}

La génération de code peut prendre plusieurs formes. Elle peut être fondée sur des templates, sur des transformations de modèles, sur des métamodèles, sur des règles déclaratives ou sur des générateurs spécifiques à un framework cible. Quelle que soit la technique choisie, l'enjeu principal n'est pas seulement de produire du code, mais de produire un code cohérent avec l'architecture attendue.

Syriani et al. (2018), dans leur cartographie des approches de génération de code par templates, montrent que les templates demeurent une technique très utilisée pour produire des artefacts logiciels répétitifs \cite{syriani2018}. Leur intérêt est évident : ils permettent de transformer une structure abstraite en fichiers concrets en conservant un contrôle relativement fin sur la forme du code. Cependant, les templates peuvent aussi devenir difficiles à maintenir lorsque les règles de génération sont dispersées, mal documentées ou trop fortement dépendantes d'un framework cible.

Greifenberg et al. (2015) analysent plusieurs mécanismes d'intégration entre code généré et code écrit manuellement dans les langages orientés objet [6]. Leur travail montre que les générateurs doivent prévoir explicitement la cohabitation entre automatisation et intervention humaine. Cette observation est particulièrement importante pour notre recherche. En effet, une application générée à partir de BPMN devra souvent être adaptée par des développeurs : ajout de règles métier spécifiques, personnalisation de l'interface, intégration avec un service externe, modification de la politique de sécurité ou optimisation d'une requête.

La transformation progressive offre une réponse à cette complexité. Au lieu de générer directement l'ensemble de l'application à partir de BPMN, il est préférable de construire une chaîne en plusieurs étapes : extraction du modèle, validation, production de contrats intermédiaires, génération du backend, génération du frontend, génération des tests et production des scripts de configuration. Cette progression rend le processus plus lisible, plus testable et plus maintenable.

\subsection{Exigences de qualité du code généré}

Le code généré doit être évalué selon des critères proches de ceux du code écrit manuellement. Il doit être lisible, structuré, cohérent avec les conventions du framework cible, facilement testable et compatible avec les outils de versionnement. Un code généré mais difficile à comprendre peut devenir un obstacle pour les développeurs. Il risque de créer une dette technique, même s'il accélère initialement le développement.

Rahad et al. (2021) étudient la place de l'humain dans la boucle de l'ingénierie dirigée par les modèles à partir de l'intégration de code écrit manuellement dans des dépôts MDE [7]. Leur contribution montre que l'intervention humaine reste importante dans les projets model-driven. Les développeurs ne disparaissent pas du processus ; ils changent plutôt de rôle. Ils doivent comprendre les modèles, adapter le code, contrôler la qualité et assurer l'évolution de l'application.

Cette observation renforce l'intérêt d'une approche où le générateur n'est pas conçu comme un substitut complet au développeur, mais comme un assistant de production logicielle. Il automatise les parties répétitives, tout en laissant aux développeurs la possibilité de comprendre, d'étendre et de corriger le code produit. Le caractère \textit{source-owned} devient ainsi une condition de qualité, et non un simple choix technique.

\mySection{Évolution, régénération et maîtrise du code généré}{}

\subsection{Le problème de la régénération}

La génération initiale d'une application constitue seulement une première étape. Dans un projet réel, les processus métier évoluent, les règles changent, les formulaires sont enrichis, de nouveaux rôles sont introduits et des intégrations externes peuvent être ajoutées. Une approche générative doit donc traiter la question de la régénération. Cette problématique est plus complexe que la génération initiale, car elle doit tenir compte des modifications déjà apportées par les développeurs.

Si la régénération écrase les fichiers modifiés manuellement, elle devient dangereuse. Si elle ne met jamais à jour le code existant, elle perd son intérêt. L'enjeu consiste donc à définir une stratégie permettant de distinguer les parties entièrement générées, les zones personnalisables, les extensions manuelles et les fichiers protégés. Cette distinction est essentielle pour que les développeurs acceptent d'utiliser un générateur dans la durée.

Dragaš et al. (2024) proposent SeamlessMDD, un framework visant l'intégration progressive du code généré et du code écrit à la main [8]. Leur approche met en avant la nécessité de préserver les pratiques existantes de développement tout en introduisant des mécanismes model-driven. Cette contribution est particulièrement utile pour notre recherche, car elle montre que la réussite d'une approche générative dépend autant de la qualité technique du générateur que de son intégration dans le flux de travail des développeurs.

\subsection{Préservation des personnalisations}

La préservation des personnalisations est une exigence centrale pour toute approche \textit{source-owned}. Lorsqu'un développeur modifie une vue, ajoute une méthode métier ou adapte une validation, ces changements ne doivent pas être perdus lors d'une nouvelle génération. Plusieurs stratégies peuvent être envisagées : séparation des fichiers générés et manuels, zones protégées, héritage, composition, génération partielle ou mécanismes de fusion.

Greifenberg et al. (2015) comparent précisément différents mécanismes d'intégration entre code généré et code manuel [6]. Leur analyse montre qu'aucune stratégie n'est universellement supérieure ; le choix dépend du langage, du framework, du niveau de personnalisation attendu et de la structure du générateur. Pour notre recherche, cette conclusion implique que l'architecture générée doit être pensée dès le départ pour séparer les responsabilités. Les fichiers entièrement régénérables doivent être distingués des fichiers destinés à l'extension humaine.

La préservation des personnalisations peut également passer par des contrats stables. Si le modèle BPMN change, le générateur peut recalculer les contrats et identifier les différences : nouvelle tâche, rôle supprimé, transition modifiée, formulaire enrichi. Cette comparaison peut ensuite guider la régénération, en signalant les conflits possibles au lieu de les résoudre silencieusement.

\subsection{Traçabilité entre modèles et fichiers générés}

La traçabilité est un autre élément important. Dans une application générée, il doit être possible de savoir quel élément du modèle BPMN a produit quel artefact logiciel. Cette relation permet de comprendre l'origine d'un fichier, de diagnostiquer une erreur, de vérifier la couverture du modèle et de faciliter l'évolution.

Dans le cas d'une application générée à partir de BPMN, la traçabilité peut être envisagée à plusieurs niveaux. Une tâche BPMN peut être liée à un formulaire frontend, à une route backend, à une méthode de service, à une règle d'autorisation et à un test. Une passerelle peut être liée à une condition de transition. Un rôle peut être lié à une politique d'accès. Ce suivi permet de transformer la génération en processus contrôlable et non en boîte noire.

La traçabilité renforce également la crédibilité académique et technique du framework. Elle permet d'évaluer la couverture fonctionnelle de la génération, de mesurer les artefacts produits et de démontrer que le passage du modèle au code repose sur des règles explicites.

\mySection{Low-code, no-code et ingénierie dirigée par les modèles}{}

\subsection{Proximité entre low-code et MDE}

Les plateformes low-code et no-code visent à réduire l'effort de développement en s'appuyant sur des modèles visuels, des composants préconfigurés, des connecteurs et des mécanismes d'automatisation. Leur objectif est de permettre la création rapide d'applications, parfois par des utilisateurs non spécialistes du développement logiciel. Cette promesse rejoint en partie celle de l'ingénierie dirigée par les modèles.

Bucaioni et al. (2022) réalisent une revue multivocale sur la modélisation dans le développement low-code [12]. Les auteurs montrent que les plateformes low-code s'appuient souvent sur des représentations visuelles ou déclaratives, même lorsqu'elles ne se réclament pas explicitement de la tradition MDE. Cette observation confirme l'existence d'un rapprochement entre les deux domaines.

Di Ruscio et al. (2022) analysent plus directement la relation entre low-code et MDE en posant la question de savoir si ces approches représentent deux faces d'une même réalité [13]. Leur analyse montre que les deux familles partagent des objectifs communs, mais diffèrent par leur culture, leurs outils, leurs utilisateurs cibles et leur rapport au code. Le low-code est souvent orienté productivité et expérience utilisateur, tandis que le MDE met davantage l'accent sur les métamodèles, les transformations, la précision formelle et la traçabilité.

Pankowska (2024) complète cette lecture en discutant le cycle de développement low-code \cite{pankowska2024}. Son analyse rappelle que la rapidité promise par le low-code ne supprime pas les étapes classiques du développement : compréhension du besoin, modélisation, configuration, intégration, test, déploiement et maintenance. Cette remarque est importante, car elle évite de présenter le low-code comme une solution magique. Même lorsque l'interface de développement est visuelle, les décisions d'architecture et de gouvernance restent présentes.

\subsection{Apports des plateformes low-code et no-code}

Les plateformes low-code/no-code apportent plusieurs avantages : accélération du développement, réduction du codage répétitif, implication des utilisateurs métier, standardisation de certains composants et facilité de prototypage. Rokis et Kirikova (2023) synthétisent ces apports dans une revue de la littérature consacrée au développement low-code [14]. Ils montrent que ces plateformes peuvent améliorer la collaboration entre équipes métier et équipes techniques, tout en réduisant le temps nécessaire à la production d'applications.

Ajimati et al. (2025) proposent une revue systématique sur l'adoption des technologies low-code/no-code dans les pratiques de développement citoyen et de transformation numérique [15]. Leur travail met en évidence le rôle croissant de ces technologies dans les organisations, mais aussi les défis associés : gouvernance, sécurité, compétences, qualité des applications, intégration avec les systèmes existants et dépendance aux plateformes.

Krouwel et al. (2024) montrent, dans un autre registre, comment des modèles d'entreprise peuvent être transformés vers une application low-code Mendix [18]. Leur travail illustre la possibilité d'utiliser une approche model-driven pour produire une application dans un environnement low-code. Cette contribution est intéressante pour notre recherche, car elle montre que les frontières entre modèle d'entreprise, modèle applicatif et plateforme d'exécution deviennent de plus en plus perméables.

\subsection{Limites du low-code pour une approche source-owned}

Malgré leurs avantages, les plateformes low-code/no-code présentent des limites importantes lorsqu'une recherche vise la production d'applications web maîtrisables et évolutives. La première limite concerne la dépendance à la plateforme. Lorsque l'application est fortement liée à un environnement propriétaire, il peut devenir difficile d'en extraire le code, de le versionner librement ou de le modifier en dehors de la plateforme. La deuxième limite concerne la profondeur de personnalisation. Certaines plateformes permettent d'aller vite pour des cas standards, mais deviennent contraignantes lorsque l'application exige une architecture spécifique.

La troisième limite concerne la gouvernance technique. Une application créée rapidement par composition visuelle peut devenir difficile à maintenir si les règles métier, les dépendances, les versions et les responsabilités ne sont pas clairement organisées. Ajimati et al. (2025) soulignent que l'adoption du low-code/no-code soulève des enjeux liés à la gouvernance, aux compétences, à la sécurité et à l'intégration organisationnelle [15].

Ces limites ne conduisent pas à rejeter le low-code. Elles montrent plutôt qu'une approche générative orientée développeurs peut reprendre certains objectifs du low-code, comme l'accélération du développement et l'abstraction, tout en évitant ses faiblesses lorsqu'elles empêchent la maîtrise du projet. L'enjeu devient alors de produire du code source exploitable, plutôt qu'une application enfermée dans une plateforme.

\mySection{Synthèse critique des travaux récents}{}

Les travaux étudiés peuvent être regroupés en cinq familles principales. La première famille concerne les approches qui utilisent BPMN comme fondement d'architectures web orientées processus. La deuxième regroupe les travaux de transformation BPMN vers des modèles d'interaction, notamment IFML. La troisième porte sur la génération d'API, de services et d'artefacts backend. La quatrième concerne l'intégration entre code généré et code manuel. La cinquième regroupe les travaux sur le low-code, le no-code et leurs liens avec le MDE.

\begin{table}[H]
\caption{Synthèse critique des familles de travaux}
\label{tab:synthese-critique-familles-travaux}
\centering
\footnotesize
\renewcommand{\arraystretch}{1.2}
\begin{tabular}{|p{2.8cm}|p{3.4cm}|p{3.5cm}|p{3.5cm}|}
\hline
\textbf{Famille de travaux} & \textbf{Apport principal} & \textbf{Limite observée} & \textbf{Enjeu pour la recherche} \\
\hline
BPMN et plateformes orientées processus & BPMN peut structurer l'exécution d'une plateforme web lorsqu'il est associé à une architecture et à un moteur de processus [2]. & L'approche reste souvent centrée sur l'orchestration plutôt que sur la génération complète du code applicatif. & Étudier comment passer d'un processus modélisé vers un projet web maîtrisable. \\
\hline
BPMN vers IFML et interfaces & Les modèles BPMN peuvent être transformés vers des modèles d'interaction et contribuer à la génération d'interfaces [3], [4]. & Les aspects backend, sécurité, tests, déploiement et régénération sont moins couverts. & Analyser les conditions d'une génération backend/frontend plus intégrée. \\
\hline
API et services autour des workflows & Les modèles de processus et de données peuvent être utilisés pour générer des API REST orientées domaine [5]. & La génération d'API ne suffit pas à produire une application complète. & Relier les contrats métier, les services et les composants d'interface. \\
\hline
Code généré et code manuel & Les approches MDE doivent prévoir l'intégration entre code généré et code écrit par les développeurs [6], [8]. & La régénération peut provoquer des conflits ou écraser des personnalisations si elle n'est pas contrôlée. & Identifier des mécanismes de séparation, de traçabilité et de préservation des personnalisations. \\
\hline
Low-code/no-code et MDE & Le low-code partage avec le MDE l'objectif d'abstraction et d'automatisation [12], [13], [15]. & Les plateformes peuvent introduire une dépendance forte, une gouvernance complexe et une faible maîtrise du code. & Conserver l'automatisation tout en préservant l'accès au code source. \\
\hline
\end{tabular}
\end{table}

Cette synthèse montre que les travaux existants apportent des réponses partielles au problème étudié. Certains clarifient le rôle de BPMN, d'autres traitent les interfaces, d'autres les API, d'autres encore la relation entre code généré et code manuel. Cependant, peu d'approches combinent dans une même chaîne la modélisation BPMN, les contrats intermédiaires, la génération backend/frontend, la maîtrise du code source et la régénération contrôlée.

Les articles complémentaires du corpus local permettent de préciser cette conclusion. Les travaux sur les applications web orientées processus montrent qu'il existe une demande réelle pour des plateformes capables de relier BPMN et application web \cite{neto2022,uyanik2023}. Les travaux sur les interfaces issues de BPMN montrent que des correspondances sont possibles, mais qu'elles nécessitent souvent des stéréotypes, des modèles d'interaction ou des informations de domaine \cite{diaz2018,laaz2020,sajji2023}. Les travaux sur la génération de code et les transformations entre artefacts soulignent enfin que la maintenabilité dépend de la clarté des transformations, de la qualité des templates et de la traçabilité \cite{syriani2018,siahaan2024}. L'état de l'art converge donc vers une même idée : la difficulté n'est pas de générer quelque chose à partir de BPMN, mais de générer un projet web complet, compréhensible et capable d'évoluer.

\mySection{Lacunes identifiées dans les travaux existants}{}

\subsection{Insuffisance de BPMN seul}
BPMN représente efficacement les processus, mais il ne décrit pas entièrement les données, les interfaces, les règles de sécurité, les validations ou les composants logiciels. Les travaux de Sajji et al. (2022, 2023) confirment cette limite, puisqu'ils complètent BPMN par IFML pour traiter les interactions utilisateur [3], [4].

Cette lacune montre que le modèle BPMN doit être considéré comme une source principale, mais non comme l'unique source d'information. Les travaux étudiés suggèrent donc la nécessité de spécifications complémentaires ou de contrats intermédiaires.

\subsection{Génération souvent partielle}

Certaines approches génèrent des interfaces ; d'autres produisent des API ; d'autres s'intéressent aux moteurs de processus ; d'autres encore étudient les modèles ou les plateformes low-code. Cependant, la construction d'une application web complète nécessite une chaîne plus intégrée.

Une application générée doit inclure au minimum une structure backend, des routes, des services, des modèles de données, des composants frontend, des formulaires, des règles d'autorisation, des tests et des scripts de configuration. Cette exigence dépasse les contributions centrées sur un seul type d'artefact et met en évidence le besoin d'approches plus globales, mais toujours modulaires.

\subsection{Maîtrise insuffisante du code source}

Les BPMS et les plateformes low-code peuvent accélérer le développement, mais ils ne garantissent pas toujours que le code source soit lisible, exportable, versionnable et modifiable. Or, pour un mémoire orienté architecture et méthodes de génération, la question de la maîtrise du code est centrale.

Greifenberg et al. (2015) et Rahad et al. (2021) montrent que l'intervention humaine reste nécessaire dans les approches model-driven [6], [7]. Dragaš et al. (2024) soulignent également l'importance d'une intégration souple entre code généré et code manuel [8]. Ces travaux justifient l'intérêt d'une orientation \textit{source-owned} : le générateur doit produire un projet logiciel que l'équipe peut réellement posséder et faire évoluer.

\subsection{Régénération et évolution peu traitées}

Dans de nombreux travaux, la génération est présentée comme une opération de production initiale. Or, les applications métier évoluent continuellement. Les modèles BPMN peuvent changer, les règles métier peuvent être modifiées, de nouveaux écrans peuvent être ajoutés et des intégrations externes peuvent apparaître.

Une approche générative complète doit donc traiter la génération comme un processus évolutif. Cela implique de détecter les différences entre versions de modèles, d'identifier les fichiers impactés, de protéger les personnalisations humaines et de signaler les conflits. La régénération contrôlée devient ainsi une dimension centrale du problème.

\subsection{Positionnement de la recherche}

À partir des lacunes identifiées, notre recherche peut être située dans le champ des approches génératives pour applications web métier pilotées par des modèles BPMN. Elle ne relève ni d'un simple éditeur BPMN, ni d'un BPMS classique, ni d'une plateforme low-code fermée. Elle se situe à l'intersection de trois familles : la modélisation BPMN, l'ingénierie dirigée par les modèles et la génération d'applications web.

L'orientation de recherche qui se dégage de cette revue repose sur cinq principes :

\begin{itemize}
    \item utiliser BPMN comme point d'entrée métier ;
    \item enrichir le modèle par des contrats intermédiaires exploitables par les générateurs ;
    \item générer des artefacts backend et frontend cohérents ;
    \item préserver la maîtrise du code source par une approche \textit{source-owned} ;
    \item prévoir l'évolution, la régénération et la modularité de l'application.
\end{itemize}

Ce positionnement permet de préparer la suite du mémoire. Après avoir identifié les limites des approches existantes, il devient possible de formuler une contribution architecturale et méthodologique répondant aux besoins dégagés par l'état de l'art.

\mySection{Conclusion}{}

À l’issue de l’analyse des travaux existants, il apparaît que les solutions actuelles — qu’il s’agisse des BPMS traditionnels, des approches basées sur l’ingénierie dirigée par les modèles ou encore des plateformes low-code — apportent chacune des réponses pertinentes à certains besoins spécifiques. Toutefois, ces solutions restent généralement limitées à des aspects particuliers du cycle de développement. Certaines mettent l’accent sur la modélisation des processus métier, d’autres sur la génération d’interfaces ou l’automatisation des workflows, sans réellement proposer une démarche unifiée permettant de passer, de manière fluide, d’un modèle BPMN à une application web complète et exploitable.

L’étude menée met également en évidence plusieurs insuffisances récurrentes. D’une part, la génération automatique reste souvent partielle et ne couvre pas l’ensemble des composants nécessaires à une application fonctionnelle. D’autre part, la cohabitation entre le code généré automatiquement et les développements manuels demeure difficile à gérer, ce qui complique la maintenance et les évolutions futures. À cela s’ajoutent la dépendance fréquente à des plateformes propriétaires ainsi qu’un faible contrôle sur le code produit, limitant ainsi sa réutilisation, sa personnalisation et sa pérennité.

Ces observations montrent donc l’absence d’un cadre global capable d’accompagner efficacement la transformation d’un processus métier modélisé en BPMN vers une application web dont le code source reste accessible, maintenable et évolutif. C’est dans cette perspective que s’inscrit la contribution de ce mémoire. Le chapitre suivant présente \textbf{ProOWeb}, un framework génératif et modulaire conçu pour répondre à ces limitations. Cette approche repose sur une modélisation BPMN enrichie, l’utilisation de contrats intermédiaires et la génération automatisée d’artefacts web, tout en garantissant une meilleure maîtrise et une évolution durable du code applicatif produit.
